% Compléments Bloc 1 — zooms de preuve pour aérer la soutenance

\begin{frame}{C1.1 — Faisabilité : croiser périmètre, coût et risque}
\context{B1}{C1.1 · Lynx Immo · Zoom preuve}
\begin{columns}[T]
\begin{column}{0.48\linewidth}
\begin{block}{Ce qu'il fallait vérifier}
\begin{tightitems}
  \item Cohérence du MVP avec le besoin client.
  \item Architecture réaliste pour l'équipe disponible.
  \item Coût d'exploitation après livraison.
  \item Fonctions à différer pour protéger la trajectoire.
\end{tightitems}
\end{block}
\end{column}
\begin{column}{0.48\linewidth}
\begin{block}{Éléments de preuve}
\begin{tightitems}
  \item Priorisation du périmètre.
  \item Comparaison simplifiée des coûts.
  \item Scénarios d'architecture.
  \item Arbitrages entre MVP et perspectives.
\end{tightitems}
\end{block}
\end{column}
\end{columns}
\begin{block}{Message jury}
La faisabilité a servi d'outil de décision et de cadrage, pas de simple réponse oui/non.
\end{block}
\end{frame}

\begin{frame}{C1.3 — Pourquoi réécrire Ediacara n'était pas la bonne réponse}
\context{B1}{C1.3 · FCU / Ediacara · Zoom preuve}
\begin{columns}[T]
\begin{column}{0.46\linewidth}
\begin{block}{La tentation initiale}
Profiter du FCU pour reconstruire l'orchestration avec des technologies modernes.
\end{block}
\begin{tightitems}
  \item Chaîne de production utilisée quotidiennement.
  \item Composants Eiffel et Scala partagés avec d'autres outils.
  \item Règles métier implicites dans le legacy.
  \item Délai et risque incompatibles avec une refonte globale.
\end{tightitems}
\end{column}
\begin{column}{0.50\linewidth}
\begin{block}{Décision retenue}
Moderniser la couche Python maîtrisable et construire des frontières sûres autour des composants trop risqués à remplacer.
\end{block}
\centering
\repofigure[0.31\textheight]{figures/fcu/tampon_executables_eiffel.png}
\end{column}
\end{columns}
\end{frame}

\begin{frame}{C1.6 — Des KPI qui doivent conduire à une décision}
\context{B1}{C1.6 · Lynx Immo · Zoom preuve}
\scriptsize
\begin{tabularx}{\linewidth}{>{\bfseries}p{3.0cm}X X}
\toprule
Famille & Ce que je mesure & Décision possible \\
\midrule
Projet & tâches, jalons, couverture des exigences & replanner, scinder, déscoper \\
Valeur / usage & engagement, fréquence, continuité & vérifier la pertinence produit \\
Technique & réponse, stabilité, dette technique & corriger ou prioriser une amélioration \\
Charge contrôlée & comportement sous sollicitation & préparer la future mise en production \\
\bottomrule
\end{tabularx}
\vspace{0.7em}
\begin{block}{Question centrale}
Qu'est-ce que cette mesure doit me faire changer dans la trajectoire du projet ?
\end{block}
\end{frame}

\begin{frame}{C1.7 — Superviser par une boucle d'arbitrage explicite}
\context{B1}{C1.7 · Lynx Immo · Zoom preuve}
\begin{center}
\begin{tikzpicture}[node distance=0.8cm,every node/.style={font=\small,align=center}]
  \node[draw=Line,rounded corners,fill=SoftBlue,minimum width=2.7cm,minimum height=0.9cm] (obs) {Observer\\l'avancement};
  \node[draw=Line,rounded corners,fill=SoftOrange,minimum width=2.7cm,minimum height=0.9cm,right=of obs] (cmp) {Comparer\\aux exigences};
  \node[draw=Line,rounded corners,fill=SoftRed,minimum width=2.7cm,minimum height=0.9cm,right=of cmp] (risk) {Qualifier\\le risque};
  \node[draw=Line,rounded corners,fill=SoftGreen,minimum width=2.7cm,minimum height=0.9cm,right=of risk] (dec) {Décider\\et réajuster};
  \draw[->,thick,Blue] (obs)--(cmp); \draw[->,thick,Blue] (cmp)--(risk); \draw[->,thick,Blue] (risk)--(dec);
\end{tikzpicture}
\end{center}
\vspace{0.8em}
\begin{tightitems}
  \item Réallouer lorsque la disponibilité change.
  \item Simplifier une fonctionnalité pour protéger le MVP.
  \item Intégrer le temps d'apprentissage dans la charge réelle.
  \item Documenter les conséquences d'un arbitrage plutôt que masquer l'écart au plan.
\end{tightitems}
\end{frame}
